\section{Introduction}

Artificial subsurface (tile) drainage is among the most consequential
anthropogenic modifications of the field-to-stream pathway in agricultural
landscapes. By lowering the water table and routing excess soil water rapidly
to outlets, drainage converts poorly drained and waterlogged land into some of
the world's most productive and intensively farmed cropland. Its footprint is
correspondingly large: roughly $12$--$17\%$ of cultivated land is artificially
drained worldwide \parencite{Valipour2014}, about three-quarters of surveyed
United States producers report drained cropland \parencite{Morton2015}, nearly
$90\%$ of Finnish farmers rate subsurface drainage as important---ahead of
irrigation \parencite{PeltonenSainio2021}, and remote-sensing inventories now
place the drained United States extent at more than $22$~Mha
\parencite{Valayamkunnath2020}. Because tile networks intercept water that would
otherwise recharge shallow groundwater or leave as surface runoff and deliver it
directly to streams, drainage behaves as a master variable of drained-land
hydrology and water quality: in humid settings tiles commonly carry a large and
measurable fraction of streamflow and dominate field water export
\parencite{King2014,Williams2018,Schilling2019}. Understanding this engineered
re-plumbing of the landscape is therefore central to managing both agricultural
productivity and the quality of the waters that receive drainage.

The hydrologic changes wrought by drainage are double-edged. Relative to
undrained or surface-drained land, improved subsurface drainage generally
reduces surface runoff, peak flow, sediment, and particulate phosphorus, yet it
opens an efficient conduit for nitrate, coupling tile-drained cropland to
hypoxia in the Gulf of Mexico and to the eutrophication of inland waters
\parencite{FAUSEY1995,SKAGGS1994,THOMAS1995,Gramlich2018,Raymond2012}. The
once-common presumption that the soil profile filters subsurface water into a
``clean'' pathway has been overturned: a small volume of hydraulically active
macropores delivers phosphorus, sediment, and conservative tracers to drains
within hours of application or rainfall \parencite{Pease2018}. At the
same time drainage remains a foundational production technology that raises
yields on excess-water-stressed soils, reclaims saline and waterlogged land, and
returns favourable financial and social benefits \parencite{Singh2022b,Yang2025}.
Reconciling these competing outcomes---water-quality protection against
agronomic and economic gain---is the defining management challenge of drained
agriculture.

The knowledge synthesised here rests on several mutually reinforcing lines of
evidence. Continuously monitored plots, fields, and watershed outlets quantify
flows and loads; tracer, isotope, and geochemical studies partition flow
pathways and fingerprint contaminant sources; controlled columns, tanks, and
lysimeters isolate mechanisms; long-term records spanning decades reveal trends
and legacies; and a mature toolkit of process-based and data-driven models,
together with near-surface geophysics and remote sensing, simulates and locates
buried networks
\parencite{Smith2011,Singh2019b,Gowda2012,Allred2020}. Where these independent
approaches converge, findings can be stated with strong confidence, and it is
that consolidated evidence---rather than open questions---that this review sets
out to establish.

This paper is the state-of-knowledge component of a two-part synthesis. It
consolidates what the evidence has established about agricultural subsurface
drainage---the consolidated findings, mechanisms, and quantified outcomes that
are well supported---and how strongly each is supported; the open questions,
contested results, and emerging frontiers are treated in a companion frontier
component. The organisation follows the causal logic of drained-land systems.
Because hydrology is the master control on everything that follows, we begin
with drainage flow pathways, the modified flow regime, and its scaling. We then
turn to contaminant loss and its mechanisms---nitrate export and its biophysical
controls; phosphorus and sediment loss and the preferential transport that
drives it; and the further dimensions of pesticides, microbes, trace elements,
and greenhouse gases. Management is treated next along the pathway from field to
outlet: in-field nutrient and cropping practices, drainage design and
water-table management, and edge-of-field treatment technologies with their
watershed-scale integration. We then consolidate the agronomic, reclamation, and
socio-economic benefits of drainage, and close with the modeling, measurement,
and mapping toolkit on which the physical and management conclusions rest.
Throughout, quantified magnitudes are developed within their respective thematic
sections; the aim of this synthesis is to bring those threads together into a
coherent account of what is now known about tile drainage.
